% Generated draft: requires mathematical and source review.
\documentclass[11pt]{article}
\usepackage[margin=25mm]{geometry}
\usepackage{fontspec}
\setmainfont{Latin Modern Roman}
\usepackage{amsmath,amssymb,mathtools,unicode-math}
\setmathfont{Latin Modern Math}
\usepackage{enumitem,needspace,xurl,hyperref}
\setlength{\parindent}{0pt}
\setlength{\parskip}{5pt}
\setlength{\emergencystretch}{3em}
\newcommand{\N}{\mathbb{N}}
\newcommand{\Z}{\mathbb{Z}}
\newcommand{\Q}{\mathbb{Q}}
\newcommand{\R}{\mathbb{R}}
\newcommand{\C}{\mathbb{C}}
\newcommand{\F}{\mathbb{F}}
\newcommand{\A}{\mathbb{A}}
\newcommand{\PP}{\mathbb{P}}
\newcommand{\E}{\mathbb{E}}
\newcommand{\eps}{\varepsilon}
\newcommand{\dd}{\,\mathrm{d}}
\begin{document}
\section*{UnsolvedMath ID: 20000981}
\textbf{Source problem code:} catalogue entry
\subsection{Definitions and mathematical statement}
Wesley Pegden's source asks whether all hexagonal plane tilings by quarter-turn-symmetric polyominoes, up to suitable dilations, are the tilings in his paper with Lionel Levine and Charles Smart [23] concerning Apollonian patterns in the sandpile scaling limit. It says that paper contains a stronger conjecture. Neither the full family nor the stronger statement is supplied, and the terms ``hexagonal'', ``regular'' and ``primitive'' require conventions.

The aid's partial result concerns regular lattice translations of one topological-disk polyomino. Represent the tile by a finite cell set $S\subset\Z^2$, whose union of closed unit squares is a topological disk and is invariant under a quarter turn about an allowed grid center. Let $\Lambda\subset\Z^2$ be a full-rank translation lattice. The reported finite tiling certificate is that $S$ is a complete coset transversal for $\Z^2/\Lambda$, equivalently that every lattice cell belongs to exactly one translate $S+\lambda$. In particular
\[
 |S|=[\Z^2:\Lambda],\qquad |S|\equiv0\text{ or }1\pmod4.
\]
A finite list labeling boundary contacts by their translation displacements determines which translates meet the tile. The actual lattice, shape and boundary labels must be supplied to check a candidate.

Shared-edge adjacency and intersection of closed tile footprints are different contact conventions; a corner contact counts only for the latter. Under the aid's six-neighbor and primitiveness conventions, its classification below area eight has the following conclusions:
\begin{itemize}
\item Area one gives the unit square case and does not yield the required six-neighbor tiling.
\item At area four, the only six-neighbor lattice case, up to the appropriate lattice symmetries, is the staggered two-by-two square tiling. It admits a domino refinement and is therefore nonprimitive in the report's intended sense.
\item At area five, the only possible quarter-turn-symmetric disk shape is the cross pentomino. Its lattice tilings have contact degree four by shared edges and eight by closed-footprint intersection, never six.
\end{itemize}
The area congruence excludes the other positive areas below eight. Consequently the aid reports that every primitive regular six-neighbor tiling in its conjectural setting has at least eight cells per tile. This is a necessary small-area obstruction, not an existence theorem at area eight or a classification of all larger tilings.

The complete definition of primitive and the exact equivalence under dilations or coordinate symmetries are not included in the snapshot. In particular, the report's inference from domino refinement to nonprimitiveness must be checked against the definition in [23]. The finite lattice certificate and small-area calculation do not settle the source's potentially broader class of all hexagonal tilings, including any nonlattice arrangements.
\subsection{Short English statement}
The aid supplies finite tests and rules out primitive regular six-neighbor tiles of area below eight. The two-by-two square case refines into dominoes, and the cross pentomino has the wrong contact count. This restricted obstruction does not identify the full family referred to in the original sandpile conjecture.
\subsection{Sources}
\begin{itemize}
\item Local UnsolvedMath registry, permanent ID $20000981$, source code \texttt{catalogue entry}, supplied complete source record and research summary.
\item AIM workshop \emph{Generalizations of chip-firing and the critical group}, Problem 16, Wesley Pegden; supplied source \url{https://aimath.org/pastworkshops/chipfiringproblems.pdf}. Canonical local record \texttt{aim-combinatorics-notes.json}, \texttt{notes[105]}.
\item Lionel Levine, Wesley Pegden and Charles Smart, source reference [23], described as a recent paper on Apollonian circle-packing patterns in the sandpile scaling limit. No full title, year or bibliographic entry, nor the stronger conjecture mentioned by Pegden, is supplied.
\item Research attempt 1: partial result; problem open at the recorded run, registry status partially solved; supplied rigor label ``proved'', novelty confidence ``low''. The machine-generated aid is not peer reviewed and requires independent verification before citation. No full authoritative citation for the small-area classification is supplied.
\end{itemize}
\end{document}
